(a) B/P11 geometry with a $60\times60\times3\,\mathrm{mm}$ oblate ellipsoid inclined at $60^\circ$ in a $120\,\mathrm{mm}$ cube; dashed lines locate the core boundaries. (b) Relative gain $100g/G_0$ for all 16 cases, with $\sigma_f=10\,\mathrm{S\,m^{-1}}$ in the main panel and $0.18\,\mathrm{S\,m^{-1}}$ in the inset. The pair $(d_-,d_+)$ gives lower and upper distances in mm, with $g_{ij}$ denoting layout $(i,j)$; blue circles/red triangles denote A/B. (c) Lower, upper and joint effects, $100(g_{17}-g_{77})/g_{77}$, $100(g_{71}-g_{77})/g_{77}$ and $100(g_{11}-g_{77})/g_{77}$, using each core/conductivity's $g_{77}$. Open/dashed and filled/solid symbols denote $0.18$ and $10\,\mathrm{S\,m^{-1}}$; the asterisk marks an effect below $0.01635\,\mu\mathrm{S}$, twice the maximum gain change among the four preselected refined states (Table S11). (d) Base- and finer-mesh joint differences for A/0.18 and B/10 on separate $\mu\mathrm{S}$ axes; $g=G_f-G_0$.
